\section{Related Work}
\label{sec:related}

\subsection{Multi-Modal Semantic Segmentation}

Multi-modal semantic segmentation aims to leverage complementary information from heterogeneous sensors to improve scene understanding beyond what any single modality can achieve.
Early approaches~\cite{hazirbas2017fusenet, sun2019rtfnet} adopted dual-encoder architectures with element-wise or concatenation-based fusion, but these methods are typically designed for fixed modal pairs (\eg, RGB-Depth or RGB-Thermal) and cannot easily generalize to arbitrary combinations.

Transformer-based architectures have significantly advanced this field.
SegFormer~\cite{xie2021segformer} introduced a hierarchical Transformer encoder with a lightweight MLP decoder, eliminating positional encoding to enable flexible multi-scale feature extraction.
Building on this backbone, CMX~\cite{zhang2022cmx} proposed a unified cross-modal fusion framework with a Cross-Modal Feature Rectification Module (CM-FRM) and a Feature Fusion Module (FFM) that performs long-range context exchange across modalities, achieving state-of-the-art results across RGB-Depth, RGB-Thermal, RGB-Polarization, and RGB-LiDAR benchmarks.
CMNeXt~\cite{zhang2023delivering} further scaled this paradigm with the Self-Query Hub (SQ-Hub) and Parallel Pooling Mixer (PPX), enabling fusion across an arbitrary number of modalities with negligible per-modality parameter overhead and demonstrating +9.10\% mIoU improvement with quad-modal input on the DELIVER benchmark.

StitchFusion~\cite{li2024stitchfusion} took a different approach by integrating large-scale pretrained models as both encoders and feature fusers, using a MultiAdapter module for cross-modal information transfer during encoding.
This enables multi-modal fusion with minimal additional parameters by leveraging pretrained representations directly.
MMSFormer~\cite{reza2024mmsformer} introduced a novel fusion block within a Transformer architecture that progressively improves performance as additional modalities are incorporated, demonstrating the effectiveness of modality-agnostic fusion on material and semantic segmentation benchmarks.

For modality robustness, Shin~\etal~\cite{shin2024crm} proposed Complementary Random Masking (CRM) for RGB-Thermal segmentation, which strategically masks input modalities during training to prevent over-reliance on any single sensor.
Combined with a self-distillation loss, CRM forces the network to learn complementary representations, improving robustness when individual modalities are degraded.
Our night augmentation pipeline draws inspiration from this approach, extending it with extreme brightness corruption and full-channel zeroing tailored for maritime nighttime scenarios.

For maritime environments specifically, the MULTIAQUA dataset~\cite{muhovic2025multiaqua} provides synchronized RGB, thermal, and LiDAR data across daytime and nighttime conditions.
It introduces robust training strategies that enable networks trained solely on daytime images to generalize to nighttime scenarios, highlighting the importance of augmentation-based domain adaptation for multi-modal maritime perception.


\subsection{Foundation Model Adaptation for Segmentation}

The Segment Anything Model (SAM)~\cite{kirillov2023sam} established a new paradigm for promptable image segmentation through large-scale pretraining on over 1 billion masks.
Its successor, SAM2~\cite{ravi2024sam2}, extended this capability to video with a streaming memory architecture that maintains temporal context through memory banks and cross-attention mechanisms.
However, both models are designed for single-modal (RGB) input and focus on instance-level segmentation, requiring adaptation for multi-modal semantic tasks.

Low-Rank Adaptation (LoRA)~\cite{hu2022lora} provides a parameter-efficient strategy for adapting large pretrained models by injecting trainable low-rank matrices into frozen attention layers.
Applied to SAM, LoRA enables task-specific adaptation while preserving the foundation model's generalization capability with minimal additional parameters.

MoE-LoRA~\cite{zhu2024moe} advanced this approach by incorporating Mixture-of-Experts (MoE)~\cite{shazeer2017moe} routing into the LoRA framework for multi-modal SAM adaptation.
By maintaining multiple expert pairs and using a gating network for top-$k$ selection, it achieves modality-adaptive feature extraction with a 32.15\% improvement under missing modality conditions.
However, the hard top-$k$ routing can lead to underutilized experts when the number of experts is small, which is common when the expert count matches the number of input modalities.
Recent work on Soft Mixtures of Experts~\cite{puigcerver2024softmoe} has shown that replacing discrete top-$k$ selection with continuous softmax weighting over all experts improves gradient flow and eliminates dead expert problems, motivating our adoption of soft gating in the LoRA context.

MemorySAM~\cite{chen2025memorysam} proposed an alternative approach by repurposing SAM2's video memory mechanism for multi-modal fusion.
By treating different modalities as sequential frames, the memory bank naturally captures cross-modal features without requiring explicit fusion modules.
Combined with a Semantic Prototype Memory Module (SPMM) that maintains category-level prototypes through exponential moving average updates, it bridges the gap from instance-level to semantic-level understanding, achieving state-of-the-art results of 65.38\% mIoU on DELIVER and 52.88\% on MCubeS.

Our work builds upon MemorySAM's memory-based multi-modal framework while introducing Soft-MoE LoRA for improved expert utilization and a hybrid modulation strategy that separates absolute quality assessment from relative importance ranking.
